%! Multiplying \& Dividing Rational Expressions — Unit 4, Lesson 4.2 — Warm-Up
\documentclass[10pt]{article}
\usepackage{algebra2-article}
\usepackage{algebra2-boxes}
\newcommand{\TallMath}[1]{$\displaystyle #1\rule[-1.4em]{0pt}{3.2em}$}

\begin{document}

\pageheader{Unit 4, Lesson 4.2}{Warm-Up}
\namedateperiod

\vspace{0.10in}

\begin{notesbox}{Getting Ready: Skills We'll Use Today}
\small Three skills, one per leg of today's lesson. Work quickly.

\vspace{4pt}
\textbf{1. Fractions --- the fast road.} Do \emph{not} multiply the big numbers out.
\begin{enumerate}[label=(\alph*), leftmargin=2.1em, itemsep=3pt]
  \item Write every part as a product of primes, divide out what matches, \emph{then} multiply:
        $\dfrac{6}{35}\cdot\dfrac{25}{9}=\dfrac{2\cdot3}{5\cdot7}\cdot\dfrac{5\cdot5}{3\cdot3}
        =\blank{1.4cm}$
  \item To \emph{divide} by a fraction, multiply by its \blank{2.4cm}:
        $\dfrac{4}{9}\div\dfrac{8}{15}=\dfrac{4}{9}\cdot\blank{1.2cm}=\blank{1.2cm}$
  \item You may never divide by \blank{1.2cm}. So in $\dfrac{4}{9}\div\dfrac{n}{15}$, the one value
        $n$ is not allowed to have is \blank{1.0cm}.
\end{enumerate}

\vspace{4pt}
\textbf{2. Factor completely} (the Unit~3 toolkit --- nothing simplifies until this is done).
\begin{enumerate}[label=(\alph*), leftmargin=2.1em, itemsep=3pt]
  \item $x^2-9=\blank{3.4cm}$
  \item $x^2+2x-8=\blank{3.4cm}$
  \item $x^2+3x=\blank{3.4cm}$ \quad (GCF)
\end{enumerate}

\vspace{4pt}
\textbf{3. Simplify and restrict} (Lesson~4.1). For $\dfrac{x^2+3x}{x^2-9}$:
\begin{enumerate}[label=(\alph*), leftmargin=2.1em, itemsep=3pt]
  \item Factored form: \blank{4.0cm} \quad Restrictions: $x\neq$ \blank{1.8cm}
  \item Simplest form: \blank{2.2cm}
  \item Which excluded value gives a \emph{hole}? \blank{1.2cm} \; Why? \blank{4.4cm}
\end{enumerate}

\end{notesbox}

\end{document}
